\section{Conclusiones}
\begin{enumerate}
  \item Los estados financieros proporcionan la base objetiva para calcular ratios.
  \item La liquidez mide capacidad de corto plazo, pero debe revisar calidad de activos corrientes.
  \item Los ratios de gestión permiten evaluar cobranza e inventarios dentro del ciclo operativo.
  \item La solvencia muestra la participación de terceros en el financiamiento.
  \item La rentabilidad evalúa la generación de utilidad respecto a ventas, activos y patrimonio.
  \item El caso integrador demuestra que los ratios se complementan y no deben leerse de forma aislada.
  \item Las decisiones empresariales deben surgir de cálculos correctos, interpretación contextual y comparación con metas.
  \item Ingeniería de Sistemas contribuye con automatización, trazabilidad, validación de datos y dashboards.
\end{enumerate}
